\documentclass[11pt,a4paper]{article}
% lesson-preamble.tex -- Shared preamble for all Phase 00 lessons
% Usage: % lesson-preamble.tex -- Shared preamble for all Phase 00 lessons
% Usage: % lesson-preamble.tex -- Shared preamble for all Phase 00 lessons
% Usage: \input{../lesson-preamble} (from subdomain directories)

% ---- Encoding and fonts ----
\usepackage[utf8]{inputenc}
\usepackage[T1]{fontenc}

% ---- Mathematics ----
\usepackage{amsmath,amssymb,amsthm}
\usepackage{mathtools}

% ---- Page geometry ----
\usepackage[margin=1in,headheight=14pt]{geometry}

% ---- Hyperlinks ----
\usepackage[colorlinks=true,linkcolor=red,citecolor=blue,urlcolor=blue]{hyperref}

% ---- Lists ----
\usepackage{enumitem}

% ---- Colors and boxes ----
\usepackage{xcolor}
\usepackage[theorems,skins,breakable]{tcolorbox}

% ---- Header/Footer ----
\usepackage{fancyhdr}
\renewcommand{\headrulewidth}{0.4pt}

% ---- Tables ----
\usepackage{booktabs}

% ---- Bibliography ----
\usepackage{natbib}
\bibliographystyle{plainnat}

% --------------------------------------------------------------------------
% Theorem environments
% --------------------------------------------------------------------------
\theoremstyle{plain}
\newtheorem{theorem}{Theorem}[section]
\newtheorem{lemma}[theorem]{Lemma}
\newtheorem{proposition}[theorem]{Proposition}
\newtheorem{corollary}[theorem]{Corollary}

\theoremstyle{definition}
\newtheorem{definition}[theorem]{Definition}
\newtheorem{example}[theorem]{Example}
\newtheorem{exercise}[theorem]{Exercise}
\newtheorem{assumption}[theorem]{Assumption}

\theoremstyle{remark}
\newtheorem{remark}[theorem]{Remark}

% --------------------------------------------------------------------------
% Colored boxes
% --------------------------------------------------------------------------
\newtcolorbox{keyresult}[1][]{%
  colback=green!5!white,
  colframe=green!50!black,
  fonttitle=\bfseries,
  title={Key Result},
  breakable,
  #1
}

\newtcolorbox{rlconnection}[1][]{%
  colback=blue!5!white,
  colframe=blue!50!black,
  fonttitle=\bfseries,
  title={RL/ML Connection},
  breakable,
  #1
}

\newtcolorbox{warning}[1][]{%
  colback=red!5!white,
  colframe=red!50!black,
  fonttitle=\bfseries,
  title={Warning},
  breakable,
  #1
}

% --------------------------------------------------------------------------
% Custom commands -- number sets
% --------------------------------------------------------------------------
\newcommand{\R}{\mathbb{R}}
\newcommand{\N}{\mathbb{N}}
\newcommand{\Q}{\mathbb{Q}}
\newcommand{\Z}{\mathbb{Z}}
\newcommand{\C}{\mathbb{C}}
\newcommand{\F}{\mathbb{F}}

% --------------------------------------------------------------------------
% Custom commands -- probability and expectation
% --------------------------------------------------------------------------
\newcommand{\E}{\mathbb{E}}
\newcommand{\Prob}{\mathbb{P}}
\newcommand{\Var}{\operatorname{Var}}
\newcommand{\Cov}{\operatorname{Cov}}
\newcommand{\indicator}{\mathbf{1}}

% --------------------------------------------------------------------------
% Custom commands -- convergence arrows
% --------------------------------------------------------------------------
\newcommand{\cas}{\xrightarrow{\text{a.s.}}}
\newcommand{\cprob}{\xrightarrow{\Prob}}
\newcommand{\clp}[1]{\xrightarrow{L^{#1}}}
\newcommand{\cdist}{\xrightarrow{d}}

% --------------------------------------------------------------------------
% Custom commands -- common operators
% --------------------------------------------------------------------------
\DeclareMathOperator{\dom}{dom}
\DeclareMathOperator{\epi}{epi}
\DeclareMathOperator{\conv}{conv}
\DeclareMathOperator{\relint}{relint}
\DeclareMathOperator{\aff}{aff}
\DeclareMathOperator{\cl}{cl}
\DeclareMathOperator{\interior}{int}
\DeclareMathOperator{\tr}{tr}
\DeclareMathOperator{\diag}{diag}
\DeclareMathOperator{\rank}{rank}
\DeclareMathOperator{\sgn}{sgn}
\DeclareMathOperator{\supp}{supp}
\DeclareMathOperator{\argmin}{arg\,min}
\DeclareMathOperator{\argmax}{arg\,max}
\DeclareMathOperator{\prox}{prox}
\DeclareMathOperator{\dist}{dist}
\DeclareMathOperator{\diam}{diam}
\DeclareMathOperator{\Span}{span}

% --------------------------------------------------------------------------
% Custom commands -- linear algebra operators
% --------------------------------------------------------------------------
\DeclareMathOperator{\nullity}{nullity}
\DeclareMathOperator{\range}{range}
\DeclareMathOperator{\nullspace}{null}
\DeclareMathOperator{\proj}{proj}

% --------------------------------------------------------------------------
% Custom commands -- measure theory
% --------------------------------------------------------------------------
\newcommand{\powerset}{\mathcal{P}}
\newcommand{\Borel}{\mathcal{B}}
 (from subdomain directories)

% ---- Encoding and fonts ----
\usepackage[utf8]{inputenc}
\usepackage[T1]{fontenc}

% ---- Mathematics ----
\usepackage{amsmath,amssymb,amsthm}
\usepackage{mathtools}

% ---- Page geometry ----
\usepackage[margin=1in,headheight=14pt]{geometry}

% ---- Hyperlinks ----
\usepackage[colorlinks=true,linkcolor=red,citecolor=blue,urlcolor=blue]{hyperref}

% ---- Lists ----
\usepackage{enumitem}

% ---- Colors and boxes ----
\usepackage{xcolor}
\usepackage[theorems,skins,breakable]{tcolorbox}

% ---- Header/Footer ----
\usepackage{fancyhdr}
\renewcommand{\headrulewidth}{0.4pt}

% ---- Tables ----
\usepackage{booktabs}

% ---- Bibliography ----
\usepackage{natbib}
\bibliographystyle{plainnat}

% --------------------------------------------------------------------------
% Theorem environments
% --------------------------------------------------------------------------
\theoremstyle{plain}
\newtheorem{theorem}{Theorem}[section]
\newtheorem{lemma}[theorem]{Lemma}
\newtheorem{proposition}[theorem]{Proposition}
\newtheorem{corollary}[theorem]{Corollary}

\theoremstyle{definition}
\newtheorem{definition}[theorem]{Definition}
\newtheorem{example}[theorem]{Example}
\newtheorem{exercise}[theorem]{Exercise}
\newtheorem{assumption}[theorem]{Assumption}

\theoremstyle{remark}
\newtheorem{remark}[theorem]{Remark}

% --------------------------------------------------------------------------
% Colored boxes
% --------------------------------------------------------------------------
\newtcolorbox{keyresult}[1][]{%
  colback=green!5!white,
  colframe=green!50!black,
  fonttitle=\bfseries,
  title={Key Result},
  breakable,
  #1
}

\newtcolorbox{rlconnection}[1][]{%
  colback=blue!5!white,
  colframe=blue!50!black,
  fonttitle=\bfseries,
  title={RL/ML Connection},
  breakable,
  #1
}

\newtcolorbox{warning}[1][]{%
  colback=red!5!white,
  colframe=red!50!black,
  fonttitle=\bfseries,
  title={Warning},
  breakable,
  #1
}

% --------------------------------------------------------------------------
% Custom commands -- number sets
% --------------------------------------------------------------------------
\newcommand{\R}{\mathbb{R}}
\newcommand{\N}{\mathbb{N}}
\newcommand{\Q}{\mathbb{Q}}
\newcommand{\Z}{\mathbb{Z}}
\newcommand{\C}{\mathbb{C}}
\newcommand{\F}{\mathbb{F}}

% --------------------------------------------------------------------------
% Custom commands -- probability and expectation
% --------------------------------------------------------------------------
\newcommand{\E}{\mathbb{E}}
\newcommand{\Prob}{\mathbb{P}}
\newcommand{\Var}{\operatorname{Var}}
\newcommand{\Cov}{\operatorname{Cov}}
\newcommand{\indicator}{\mathbf{1}}

% --------------------------------------------------------------------------
% Custom commands -- convergence arrows
% --------------------------------------------------------------------------
\newcommand{\cas}{\xrightarrow{\text{a.s.}}}
\newcommand{\cprob}{\xrightarrow{\Prob}}
\newcommand{\clp}[1]{\xrightarrow{L^{#1}}}
\newcommand{\cdist}{\xrightarrow{d}}

% --------------------------------------------------------------------------
% Custom commands -- common operators
% --------------------------------------------------------------------------
\DeclareMathOperator{\dom}{dom}
\DeclareMathOperator{\epi}{epi}
\DeclareMathOperator{\conv}{conv}
\DeclareMathOperator{\relint}{relint}
\DeclareMathOperator{\aff}{aff}
\DeclareMathOperator{\cl}{cl}
\DeclareMathOperator{\interior}{int}
\DeclareMathOperator{\tr}{tr}
\DeclareMathOperator{\diag}{diag}
\DeclareMathOperator{\rank}{rank}
\DeclareMathOperator{\sgn}{sgn}
\DeclareMathOperator{\supp}{supp}
\DeclareMathOperator{\argmin}{arg\,min}
\DeclareMathOperator{\argmax}{arg\,max}
\DeclareMathOperator{\prox}{prox}
\DeclareMathOperator{\dist}{dist}
\DeclareMathOperator{\diam}{diam}
\DeclareMathOperator{\Span}{span}

% --------------------------------------------------------------------------
% Custom commands -- linear algebra operators
% --------------------------------------------------------------------------
\DeclareMathOperator{\nullity}{nullity}
\DeclareMathOperator{\range}{range}
\DeclareMathOperator{\nullspace}{null}
\DeclareMathOperator{\proj}{proj}

% --------------------------------------------------------------------------
% Custom commands -- measure theory
% --------------------------------------------------------------------------
\newcommand{\powerset}{\mathcal{P}}
\newcommand{\Borel}{\mathcal{B}}
 (from subdomain directories)

% ---- Encoding and fonts ----
\usepackage[utf8]{inputenc}
\usepackage[T1]{fontenc}

% ---- Mathematics ----
\usepackage{amsmath,amssymb,amsthm}
\usepackage{mathtools}

% ---- Page geometry ----
\usepackage[margin=1in,headheight=14pt]{geometry}

% ---- Hyperlinks ----
\usepackage[colorlinks=true,linkcolor=red,citecolor=blue,urlcolor=blue]{hyperref}

% ---- Lists ----
\usepackage{enumitem}

% ---- Colors and boxes ----
\usepackage{xcolor}
\usepackage[theorems,skins,breakable]{tcolorbox}

% ---- Header/Footer ----
\usepackage{fancyhdr}
\renewcommand{\headrulewidth}{0.4pt}

% ---- Tables ----
\usepackage{booktabs}

% ---- Bibliography ----
\usepackage{natbib}
\bibliographystyle{plainnat}

% --------------------------------------------------------------------------
% Theorem environments
% --------------------------------------------------------------------------
\theoremstyle{plain}
\newtheorem{theorem}{Theorem}[section]
\newtheorem{lemma}[theorem]{Lemma}
\newtheorem{proposition}[theorem]{Proposition}
\newtheorem{corollary}[theorem]{Corollary}

\theoremstyle{definition}
\newtheorem{definition}[theorem]{Definition}
\newtheorem{example}[theorem]{Example}
\newtheorem{exercise}[theorem]{Exercise}
\newtheorem{assumption}[theorem]{Assumption}

\theoremstyle{remark}
\newtheorem{remark}[theorem]{Remark}

% --------------------------------------------------------------------------
% Colored boxes
% --------------------------------------------------------------------------
\newtcolorbox{keyresult}[1][]{%
  colback=green!5!white,
  colframe=green!50!black,
  fonttitle=\bfseries,
  title={Key Result},
  breakable,
  #1
}

\newtcolorbox{rlconnection}[1][]{%
  colback=blue!5!white,
  colframe=blue!50!black,
  fonttitle=\bfseries,
  title={RL/ML Connection},
  breakable,
  #1
}

\newtcolorbox{warning}[1][]{%
  colback=red!5!white,
  colframe=red!50!black,
  fonttitle=\bfseries,
  title={Warning},
  breakable,
  #1
}

% --------------------------------------------------------------------------
% Custom commands -- number sets
% --------------------------------------------------------------------------
\newcommand{\R}{\mathbb{R}}
\newcommand{\N}{\mathbb{N}}
\newcommand{\Q}{\mathbb{Q}}
\newcommand{\Z}{\mathbb{Z}}
\newcommand{\C}{\mathbb{C}}
\newcommand{\F}{\mathbb{F}}

% --------------------------------------------------------------------------
% Custom commands -- probability and expectation
% --------------------------------------------------------------------------
\newcommand{\E}{\mathbb{E}}
\newcommand{\Prob}{\mathbb{P}}
\newcommand{\Var}{\operatorname{Var}}
\newcommand{\Cov}{\operatorname{Cov}}
\newcommand{\indicator}{\mathbf{1}}

% --------------------------------------------------------------------------
% Custom commands -- convergence arrows
% --------------------------------------------------------------------------
\newcommand{\cas}{\xrightarrow{\text{a.s.}}}
\newcommand{\cprob}{\xrightarrow{\Prob}}
\newcommand{\clp}[1]{\xrightarrow{L^{#1}}}
\newcommand{\cdist}{\xrightarrow{d}}

% --------------------------------------------------------------------------
% Custom commands -- common operators
% --------------------------------------------------------------------------
\DeclareMathOperator{\dom}{dom}
\DeclareMathOperator{\epi}{epi}
\DeclareMathOperator{\conv}{conv}
\DeclareMathOperator{\relint}{relint}
\DeclareMathOperator{\aff}{aff}
\DeclareMathOperator{\cl}{cl}
\DeclareMathOperator{\interior}{int}
\DeclareMathOperator{\tr}{tr}
\DeclareMathOperator{\diag}{diag}
\DeclareMathOperator{\rank}{rank}
\DeclareMathOperator{\sgn}{sgn}
\DeclareMathOperator{\supp}{supp}
\DeclareMathOperator{\argmin}{arg\,min}
\DeclareMathOperator{\argmax}{arg\,max}
\DeclareMathOperator{\prox}{prox}
\DeclareMathOperator{\dist}{dist}
\DeclareMathOperator{\diam}{diam}
\DeclareMathOperator{\Span}{span}

% --------------------------------------------------------------------------
% Custom commands -- linear algebra operators
% --------------------------------------------------------------------------
\DeclareMathOperator{\nullity}{nullity}
\DeclareMathOperator{\range}{range}
\DeclareMathOperator{\nullspace}{null}
\DeclareMathOperator{\proj}{proj}

% --------------------------------------------------------------------------
% Custom commands -- measure theory
% --------------------------------------------------------------------------
\newcommand{\powerset}{\mathcal{P}}
\newcommand{\Borel}{\mathcal{B}}

\pagestyle{fancy}
\fancyhf{}
\fancyhead[L]{\textsc{Lesson 2: Topology and Continuity}}
\fancyhead[R]{\thepage}
\title{Lesson 2: Topology and Continuity}
\author{AI/ML Mathematics}
\date{}
\begin{document}
\maketitle
\tableofcontents
\newpage

%% ============================================================
%% SECTION 1: MOTIVATION AND OVERVIEW
%% ============================================================
\section{Motivation and Overview}
\label{sec:motivation}

Consider training a neural network $f_\theta \colon \R^d \to \R$ to minimize
a loss function $L(\theta)$ over a parameter space $\Theta \subseteq \R^d$.
A natural question is: does the minimum of $L$ over $\Theta$ actually
exist, or could the infimum be a value that is never attained? In finite
dimensions, the \emph{extreme value theorem} guarantees that a continuous
function on a \emph{compact} set attains its maximum and minimum. But what
exactly is compactness, and why does it work? Understanding this requires
the topology of the real line---open sets, closed sets, and the
Heine--Borel characterization of compact subsets of $\R^n$---together
with the precise $\varepsilon$-$\delta$ definition of continuity. The tools
developed in this lesson will let us prove that continuous functions on
closed bounded sets attain their extreme values, that the intermediate
value theorem provides existence of roots and fixed points, and that
uniform continuity gives the quantitative control needed for approximation
arguments throughout optimization and reinforcement learning.

\begin{rlconnection}[title={Why RL/ML needs topology and continuity}]
\begin{itemize}[nosep]
  \item \textbf{Existence of optimal solutions.} The extreme value theorem
    guarantees that a continuous loss function attains its minimum on a
    compact parameter set---a foundational assumption in optimization.
  \item \textbf{Value function continuity.} In MDPs with continuous state
    spaces, the value function $V^\pi$ is often shown to be continuous;
    this enables function approximation and ensures nearby states have
    similar values.
  \item \textbf{Compactness in function approximation.} The
    Arzel\`{a}--Ascoli theorem (a compactness criterion for function
    spaces) underpins universal approximation results for neural networks.
  \item \textbf{Fixed-point existence.} Brouwer's fixed-point theorem
    (every continuous map from a compact convex set to itself has a fixed
    point) is used to prove existence of Nash equilibria and optimal
    policies.
\end{itemize}
\end{rlconnection}

\paragraph{Prerequisites.}
Completeness of $\R$ (the least upper bound property), convergence of
sequences (the $\varepsilon$-$N$ definition), the algebraic limit theorems,
the Bolzano--Weierstrass theorem (Lesson~1, Theorem~4.7), and the Cauchy
criterion for convergence (Lesson~1, Theorem~4.10).

\paragraph{Learning objectives.}
After completing this lesson, you will be able to:
\begin{itemize}[nosep]
  \item \textbf{Define} open sets, closed sets, closure, interior, and
    boundary in $\R$, and \textbf{classify} sets using these concepts.
  \item \textbf{State} and \textbf{prove} the Heine--Borel theorem
    characterizing compact subsets of $\R$.
  \item \textbf{Prove} that sequential compactness and open-cover
    compactness are equivalent in $\R$.
  \item \textbf{Define} continuity via $\varepsilon$-$\delta$ and via
    sequences, and \textbf{prove} their equivalence.
  \item \textbf{Derive} the extreme value theorem and the intermediate
    value theorem from compactness and continuity.
  \item \textbf{Prove} that a continuous function on a compact set is
    uniformly continuous.
\end{itemize}

%% ============================================================
%% SECTION 2: REFERENCES
%% ============================================================
\section{References}
\label{sec:references}

\begin{itemize}
  \item \citet[Chapters 2 and 4]{rudin1976} --- Topology of the real line,
    compactness, and continuity with complete proofs.
  \item \citet[Chapters 3--4]{abbott2015} --- An accessible treatment of
    open and closed sets, compact sets, and the extreme value theorem.
  \item \citet[Chapter 4]{bartle2011} --- Detailed development of
    continuity with many worked examples and exercises.
  \item \citet[Chapters 2--3]{pugh2015} --- A geometric perspective on
    topology and continuity in $\R$ and $\R^n$.
  \item \citet[Chapter 1]{folland1999} --- Point-set topology foundations
    as preparation for measure theory.
\end{itemize}

%% ============================================================
%% SECTION 3: OPEN AND CLOSED SETS
%% ============================================================
\section{Open and Closed Sets}
\label{sec:open_closed}

\begin{definition}[$\varepsilon$-Neighborhood]
\label{def:neighborhood}
For $x \in \R$ and $\varepsilon > 0$, the \emph{$\varepsilon$-neighborhood}
of $x$ is the open interval
$V_\varepsilon(x) = (x - \varepsilon, x + \varepsilon) = \{y \in \R :
|y - x| < \varepsilon\}$.
\end{definition}

\begin{definition}[Open Set]
\label{def:open}
A set $O \subseteq \R$ is \emph{open} if for every $x \in O$ there exists
$\varepsilon > 0$ such that $V_\varepsilon(x) \subseteq O$.
\end{definition}

\begin{example}
\label{ex:open_interval}
Every open interval $(a, b)$ is an open set. If $x \in (a, b)$, take
$\varepsilon = \min(x - a, b - x) > 0$. Then
$V_\varepsilon(x) \subseteq (a, b)$.
\end{example}

\begin{example}
\label{ex:R_open}
The entire real line $\R$ is open (take any $\varepsilon > 0$), and the
empty set $\emptyset$ is open (the condition holds vacuously).
\end{example}

\begin{theorem}[Properties of Open Sets]
\label{thm:open_props}
\begin{enumerate}[nosep,label=(\roman*)]
  \item The union of any collection of open sets is open.
  \item The intersection of finitely many open sets is open.
\end{enumerate}
\end{theorem}

\begin{proof}
(i) Let $\{O_\alpha\}_{\alpha \in A}$ be a collection of open sets and
$O = \bigcup_{\alpha \in A} O_\alpha$. If $x \in O$, then $x \in O_\beta$
for some $\beta \in A$. Since $O_\beta$ is open, there exists
$\varepsilon > 0$ with $V_\varepsilon(x) \subseteq O_\beta \subseteq O$.

(ii) Let $O_1, \ldots, O_n$ be open and $x \in O_1 \cap \cdots \cap O_n$.
For each $k$, there exists $\varepsilon_k > 0$ with
$V_{\varepsilon_k}(x) \subseteq O_k$. Set
$\varepsilon = \min(\varepsilon_1, \ldots, \varepsilon_n) > 0$. Then
$V_\varepsilon(x) \subseteq O_k$ for all $k$, so
$V_\varepsilon(x) \subseteq O_1 \cap \cdots \cap O_n$.
\end{proof}

\begin{warning}[title={Infinite intersections of open sets need not be open}]
The finite intersection condition in Theorem~\ref{thm:open_props}(ii)
cannot be extended to infinite intersections.  For example,
$\bigcap_{n=1}^\infty (-1/n, 1/n) = \{0\}$, which is not open.
\end{warning}

\begin{definition}[Closed Set]
\label{def:closed}
A set $F \subseteq \R$ is \emph{closed} if its complement
$F^c = \R \setminus F$ is open.
\end{definition}

\begin{example}
\label{ex:closed_interval}
The closed interval $[a, b]$ is closed: its complement
$(-\infty, a) \cup (b, \infty)$ is a union of open sets, hence open by
Theorem~\ref{thm:open_props}(i).
\end{example}

\begin{theorem}[Sequential Characterization of Closed Sets]
\label{thm:closed_seq}
A set $F \subseteq \R$ is closed if and only if every convergent sequence
$(x_n)$ with $x_n \in F$ for all $n$ has its limit in $F$.
\end{theorem}

\begin{proof}
$(\Rightarrow)$ Suppose $F$ is closed and $x_n \to L$ with $x_n \in F$.
Assume for contradiction that $L \notin F$. Then $L \in F^c$, which is
open, so there exists $\varepsilon > 0$ with
$V_\varepsilon(L) \subseteq F^c$. But $x_n \to L$ means there exists $N$
with $|x_n - L| < \varepsilon$ for $n \ge N$, giving $x_N \in F^c$, a
contradiction since $x_N \in F$.

$(\Leftarrow)$ Suppose $F$ satisfies the sequential condition. We show
$F^c$ is open. Let $x \in F^c$. If no $\varepsilon$-neighborhood of $x$
is contained in $F^c$, then for each $n \in \N$ there exists
$x_n \in F \cap V_{1/n}(x)$. Then $x_n \to x$ with $x_n \in F$, so
$x \in F$ by hypothesis, contradicting $x \in F^c$. Hence some
$V_\varepsilon(x) \subseteq F^c$, and $F^c$ is open.
\end{proof}

\begin{definition}[Interior, Closure, Boundary]
\label{def:int_cl_bd}
Let $A \subseteq \R$.
\begin{enumerate}[nosep,label=(\roman*)]
  \item The \emph{interior} of $A$ is
    $\interior(A) = \{x \in A : \exists\, \varepsilon > 0,\;
    V_\varepsilon(x) \subseteq A\}$. It is the largest open set contained
    in $A$.
  \item The \emph{closure} of $A$ is
    $\cl(A) = \{x \in \R : \text{every } V_\varepsilon(x)
    \text{ intersects } A\}$. Equivalently, $\cl(A)$ is the smallest
    closed set containing $A$.
  \item The \emph{boundary} of $A$ is
    $\partial A = \cl(A) \setminus \interior(A)$.
\end{enumerate}
\end{definition}

\begin{proposition}[Sequential Characterization of Closure]
\label{prop:closure_seq}
$x \in \cl(A)$ if and only if there exists a sequence $(a_n)$ in $A$
with $a_n \to x$.
\end{proposition}

\begin{proof}
$(\Rightarrow)$ If $x \in \cl(A)$, then for each $n \in \N$,
$V_{1/n}(x) \cap A \neq \emptyset$, so we can choose
$a_n \in V_{1/n}(x) \cap A$. Then $|a_n - x| < 1/n$, so $a_n \to x$.

$(\Leftarrow)$ If $a_n \to x$ with $a_n \in A$, then for any
$\varepsilon > 0$, there exists $N$ with $a_N \in V_\varepsilon(x) \cap A$.
Hence $V_\varepsilon(x) \cap A \neq \emptyset$ for every $\varepsilon > 0$,
so $x \in \cl(A)$.
\end{proof}

\begin{example}
\label{ex:closure_Q}
The closure $\cl(\Q) = \R$, since every real number is the limit of a
sequence of rationals (by the density of $\Q$ in $\R$, Lesson~1,
Theorem~3.11). Similarly, $\interior(\Q) = \emptyset$, since no open
interval is contained in $\Q$ (every interval contains irrationals).
Hence $\partial \Q = \R$.
\end{example}

With the basic notions of open and closed sets in place, we now turn to
compactness, the property that makes topology most useful for analysis.

%% ============================================================
%% SECTION 4: COMPACTNESS
%% ============================================================
\section{Compactness}
\label{sec:compactness}

\begin{definition}[Open Cover, Compactness]
\label{def:compact}
Let $K \subseteq \R$.
\begin{enumerate}[nosep,label=(\roman*)]
  \item An \emph{open cover} of $K$ is a collection
    $\{O_\alpha\}_{\alpha \in A}$ of open sets whose union contains $K$:
    $K \subseteq \bigcup_{\alpha \in A} O_\alpha$.
  \item A \emph{finite subcover} is a finite subcollection
    $\{O_{\alpha_1}, \ldots, O_{\alpha_n}\}$ that still covers $K$.
  \item $K$ is \emph{compact} if every open cover of $K$ has a finite
    subcover.
\end{enumerate}
\end{definition}

\begin{definition}[Sequential Compactness]
\label{def:seq_compact}
A set $K \subseteq \R$ is \emph{sequentially compact} if every sequence
$(x_n)$ in $K$ has a subsequence that converges to a point in $K$.
\end{definition}

\begin{theorem}[Characterization of Compact Sets in $\R$]
\label{thm:compact_equiv}
For $K \subseteq \R$, the following are equivalent:
\begin{enumerate}[nosep,label=(\roman*)]
  \item $K$ is compact (every open cover has a finite subcover).
  \item $K$ is sequentially compact.
  \item $K$ is closed and bounded.
\end{enumerate}
\end{theorem}

We prove this theorem through a sequence of implications.

\begin{lemma}
\label{lem:compact_closed_bounded}
If $K$ is compact, then $K$ is closed and bounded.
\end{lemma}

\begin{proof}
\emph{Bounded:} The collection $\{(-n, n) : n \in \N\}$ is an open cover
of $K$ (since $\R = \bigcup_{n=1}^\infty (-n, n)$). By compactness,
finitely many suffice: $K \subseteq (-n_1, n_1) \cup \cdots \cup
(-n_m, n_m) = (-N, N)$ where $N = \max(n_1, \ldots, n_m)$. Hence $K$ is
bounded.

\emph{Closed:} We show $K^c$ is open. Let $y \in K^c$. For each
$x \in K$, let $\varepsilon_x = |x - y|/2 > 0$ (since $y \neq x$). The
open intervals $V_{\varepsilon_x}(x)$ cover $K$. By compactness, finitely
many suffice: $K \subseteq V_{\varepsilon_{x_1}}(x_1) \cup \cdots \cup
V_{\varepsilon_{x_m}}(x_m)$. Set
$\delta = \min(\varepsilon_{x_1}, \ldots, \varepsilon_{x_m}) > 0$. Then
$V_\delta(y) \cap V_{\varepsilon_{x_k}}(x_k) = \emptyset$ for each $k$
(since $|y - x_k| = 2\varepsilon_{x_k} \ge 2\delta$, and if $z$ were in
both then $|x_k - y| \le |x_k - z| + |z - y| < \varepsilon_{x_k} + \delta
\le 2\varepsilon_{x_k} = |x_k - y|$, a contradiction). Hence
$V_\delta(y) \subseteq K^c$, so $K^c$ is open.
\end{proof}

\begin{lemma}
\label{lem:closed_bounded_seq}
If $K$ is closed and bounded, then $K$ is sequentially compact.
\end{lemma}

\begin{proof}
Let $(x_n)$ be a sequence in $K$. Since $K$ is bounded, $(x_n)$ is
bounded. By the Bolzano--Weierstrass theorem (Lesson~1, Theorem~4.7),
$(x_n)$ has a convergent subsequence $x_{n_k} \to L$. Since $K$ is
closed, $L \in K$ by Theorem~\ref{thm:closed_seq}. Hence $K$ is
sequentially compact.
\end{proof}

\begin{lemma}
\label{lem:seq_compact_implies_compact}
If $K$ is sequentially compact, then $K$ is compact.
\end{lemma}

\begin{proof}
Let $\{O_\alpha\}_{\alpha \in A}$ be an open cover of $K$. We proceed in
two steps.

\textbf{Step 1 (Lebesgue number).} We claim there exists $\delta > 0$
such that for every $x \in K$, the interval $(x - \delta, x + \delta)$ is
contained in some $O_\alpha$. Suppose not: for each $n \in \N$, there
exists $x_n \in K$ such that $(x_n - 1/n, x_n + 1/n)$ is not contained
in any single $O_\alpha$. By sequential compactness, $(x_n)$ has a
subsequence $x_{n_k} \to L \in K$. Since $L \in K$ and
$\{O_\alpha\}$ covers $K$, there exists $\beta \in A$ with
$L \in O_\beta$. Since $O_\beta$ is open, there exists $\varepsilon > 0$
with $(L - \varepsilon, L + \varepsilon) \subseteq O_\beta$. Choose $k$
large enough that $|x_{n_k} - L| < \varepsilon/2$ and $1/n_k <
\varepsilon/2$. Then
$(x_{n_k} - 1/n_k, x_{n_k} + 1/n_k) \subseteq
(L - \varepsilon, L + \varepsilon) \subseteq O_\beta$,
contradicting the choice of $x_{n_k}$.

\textbf{Step 2 (Finite subcover).} Since $K$ is sequentially compact, it
is bounded (otherwise, choosing $x_n \in K$ with $|x_n| > n$ would give
a sequence with no convergent subsequence). Say $K \subseteq [-M, M]$.
Partition $[-M, M]$ into subintervals of length less than $\delta$; this
requires at most $N = \lceil 2M/\delta \rceil$ subintervals. Each
subinterval that intersects $K$ has diameter less than $\delta$, so it
fits inside some $O_\alpha$ (by Step~1, applied to any point of $K$ in
that subinterval). Selecting one such $O_\alpha$ for each of the at most
$N$ subintervals gives a finite subcover.
\end{proof}

\begin{proof}[Proof of Theorem~\ref{thm:compact_equiv}]
Lemma~\ref{lem:compact_closed_bounded} gives (i)$\Rightarrow$(iii),
Lemma~\ref{lem:closed_bounded_seq} gives (iii)$\Rightarrow$(ii), and
Lemma~\ref{lem:seq_compact_implies_compact} gives (ii)$\Rightarrow$(i).
\end{proof}

\begin{keyresult}[title={Heine--Borel Theorem}]
\begin{corollary}[Heine--Borel]
\label{cor:heine_borel}
A subset of $\R$ is compact if and only if it is closed and bounded.
\end{corollary}
This is the equivalence (i)$\Leftrightarrow$(iii) of
Theorem~\ref{thm:compact_equiv}. In $\R^n$, the same characterization
holds: $K \subseteq \R^n$ is compact if and only if $K$ is closed and
bounded.
\end{keyresult}

\begin{rlconnection}[title={Compactness in optimization and function approximation}]
Compactness ensures that optimization problems have solutions. When we
write $\min_{\theta \in \Theta} L(\theta)$ with $\Theta$ compact and $L$
continuous, the extreme value theorem (proved below) guarantees the minimum
is attained. In RL, compact action spaces ensure that
$\max_a Q(s,a)$ is achieved; without compactness, the supremum might not be
attained, and greedy policies might not exist. In function approximation,
the Arzel\`{a}--Ascoli theorem uses compactness (of the domain and of
function families) to extract uniformly convergent subsequences from bounded
equicontinuous families.
\end{rlconnection}

\begin{example}
\label{ex:compact_sets}
\begin{enumerate}[nosep,label=(\alph*)]
  \item $[0, 1]$ is compact (closed and bounded).
  \item $(0, 1)$ is not compact (not closed: $1/n \to 0 \notin (0,1)$).
  \item $[0, \infty)$ is not compact (not bounded).
  \item $\{1/n : n \in \N\} \cup \{0\}$ is compact (closed and bounded).
  \item $\{1/n : n \in \N\}$ is not compact (not closed: $1/n \to 0$
    but $0$ is not in the set).
\end{enumerate}
\end{example}

\begin{proposition}[Closed Subsets of Compact Sets]
\label{prop:closed_subset_compact}
Every closed subset of a compact set is compact.
\end{proposition}

\begin{proof}
Let $F \subseteq K$ with $F$ closed and $K$ compact. Then $K$ is closed
and bounded (Heine--Borel). Since $F \subseteq K$, $F$ is bounded. Since
$F$ is closed and bounded, $F$ is compact by Heine--Borel.
\end{proof}

With compactness established, we now develop continuity and show how
these two concepts interact to produce the central theorems of analysis.

%% ============================================================
%% SECTION 5: CONTINUITY
%% ============================================================
\section{Continuity}
\label{sec:continuity}

\begin{definition}[Continuity at a Point]
\label{def:continuity}
Let $f \colon A \to \R$ with $A \subseteq \R$, and let $c \in A$. Then
$f$ is \emph{continuous at $c$} if for every $\varepsilon > 0$ there
exists $\delta > 0$ such that $|f(x) - f(c)| < \varepsilon$ whenever
$x \in A$ and $|x - c| < \delta$. We say $f$ is \emph{continuous on $A$}
if it is continuous at every point of $A$.
\end{definition}

\begin{theorem}[Sequential Characterization of Continuity]
\label{thm:seq_cont}
Let $f \colon A \to \R$ and $c \in A$. Then $f$ is continuous at $c$ if
and only if for every sequence $(x_n)$ in $A$ with $x_n \to c$, we have
$f(x_n) \to f(c)$.
\end{theorem}

\begin{proof}
$(\Rightarrow)$ Suppose $f$ is continuous at $c$ and $x_n \to c$. Given
$\varepsilon > 0$, choose $\delta > 0$ from the definition of continuity.
Since $x_n \to c$, there exists $N$ with $|x_n - c| < \delta$ for
$n \ge N$. Then $|f(x_n) - f(c)| < \varepsilon$ for $n \ge N$, so
$f(x_n) \to f(c)$.

$(\Leftarrow)$ Suppose $f$ is not continuous at $c$. Then there exists
$\varepsilon_0 > 0$ such that for every $\delta > 0$ there exists
$x \in A$ with $|x - c| < \delta$ but $|f(x) - f(c)| \ge \varepsilon_0$.
Taking $\delta = 1/n$ for each $n \in \N$ yields a sequence $(x_n)$ in
$A$ with $x_n \to c$ but $|f(x_n) - f(c)| \ge \varepsilon_0$ for all
$n$, so $f(x_n) \not\to f(c)$.
\end{proof}

\begin{example}
\label{ex:cont_polynomial}
Every polynomial $p(x) = a_0 + a_1 x + \cdots + a_d x^d$ is continuous
on $\R$. By the algebraic limit theorems (Lesson~1, Theorem~4.4), if
$x_n \to c$ then $x_n^k \to c^k$ for each $k$ (by induction on $k$,
using the product rule), and
$p(x_n) = a_0 + a_1 x_n + \cdots + a_d x_n^d \to p(c)$ (by the sum
rule). The sequential characterization
(Theorem~\ref{thm:seq_cont}) then gives continuity.
\end{example}

\begin{theorem}[Algebraic Properties of Continuous Functions]
\label{thm:cont_algebra}
Let $f, g \colon A \to \R$ be continuous at $c \in A$. Then:
\begin{enumerate}[nosep,label=(\roman*)]
  \item $f + g$ is continuous at $c$.
  \item $f \cdot g$ is continuous at $c$.
  \item $f / g$ is continuous at $c$, provided $g(c) \neq 0$.
\end{enumerate}
\end{theorem}

\begin{proof}
All three follow immediately from Theorem~\ref{thm:seq_cont} and the
algebraic limit theorems (Lesson~1, Theorem~4.4): if $x_n \to c$, then
$f(x_n) \to f(c)$ and $g(x_n) \to g(c)$, so
$(f + g)(x_n) = f(x_n) + g(x_n) \to f(c) + g(c)$, and similarly for
products and quotients.
\end{proof}

\begin{theorem}[Composition of Continuous Functions]
\label{thm:cont_composition}
Let $f \colon A \to \R$ be continuous at $c$, and let $g \colon B \to \R$
be continuous at $f(c)$, where $f(A) \subseteq B$. Then $g \circ f$ is
continuous at $c$.
\end{theorem}

\begin{proof}
Let $x_n \to c$ with $x_n \in A$. Continuity of $f$ at $c$ gives
$f(x_n) \to f(c)$. Since $f(x_n) \in B$ and $g$ is continuous at $f(c)$,
we get $g(f(x_n)) \to g(f(c))$. By the sequential characterization,
$g \circ f$ is continuous at $c$.
\end{proof}

\begin{example}
\label{ex:cont_composition}
The function $h(x) = \sqrt{1 + x^2}$ is continuous on $\R$: it is the
composition of $f(x) = 1 + x^2$ (continuous, being a polynomial) and
$g(y) = \sqrt{y}$ (continuous on $[0, \infty)$).
\end{example}

\begin{theorem}[Preservation of Closed and Open Sets]
\label{thm:cont_preimage}
Let $f \colon A \to \R$ be continuous on $A$.
\begin{enumerate}[nosep,label=(\roman*)]
  \item If $O \subseteq \R$ is open, then
    $f^{-1}(O) = \{x \in A : f(x) \in O\}$ is open relative to $A$.
  \item If $F \subseteq \R$ is closed, then $f^{-1}(F)$ is closed
    relative to $A$.
\end{enumerate}
\end{theorem}

\begin{proof}
(i) Let $c \in f^{-1}(O)$, so $f(c) \in O$. Since $O$ is open, there
exists $\varepsilon > 0$ with $(f(c) - \varepsilon,
f(c) + \varepsilon) \subseteq O$. By continuity, there exists $\delta > 0$
with $|f(x) - f(c)| < \varepsilon$ whenever $x \in A$ and
$|x - c| < \delta$. Hence $\{x \in A : |x - c| < \delta\} \subseteq
f^{-1}(O)$.

(ii) $f^{-1}(F) = A \setminus f^{-1}(F^c)$, and $F^c$ is open, so
$f^{-1}(F^c)$ is open relative to $A$ by~(i). Hence $f^{-1}(F)$ is closed
relative to $A$.
\end{proof}

We now combine continuity with compactness to obtain the most powerful
results of this lesson.

%% ============================================================
%% SECTION 6: EXTREME VALUE THEOREM AND IVT
%% ============================================================
\section{Extreme Value Theorem and Intermediate Value Theorem}
\label{sec:evt_ivt}

\begin{theorem}[Image of Compact Set is Compact]
\label{thm:cont_compact}
Let $f \colon K \to \R$ be continuous and $K$ compact. Then $f(K)$ is
compact.
\end{theorem}

\begin{proof}
We use sequential compactness. Let $(y_n)$ be a sequence in $f(K)$.
For each $n$, choose $x_n \in K$ with $f(x_n) = y_n$. Since $K$ is
compact (hence sequentially compact), there exists a subsequence
$x_{n_k} \to x \in K$. By continuity, $y_{n_k} = f(x_{n_k}) \to f(x)
\in f(K)$. Hence $f(K)$ is sequentially compact, therefore compact.
\end{proof}

\begin{keyresult}[title={Extreme Value Theorem and Intermediate Value Theorem}]
\begin{theorem}[Extreme Value Theorem (EVT)]
\label{thm:evt}
If $f \colon [a, b] \to \R$ is continuous, then $f$ attains its maximum
and minimum: there exist $x_M, x_m \in [a, b]$ with
\[
  f(x_m) \le f(x) \le f(x_M) \quad \text{for all } x \in [a, b].
\]
\end{theorem}

\begin{theorem}[Intermediate Value Theorem (IVT)]
\label{thm:ivt}
If $f \colon [a, b] \to \R$ is continuous and $L$ is any value between
$f(a)$ and $f(b)$ (i.e., $f(a) \le L \le f(b)$ or
$f(b) \le L \le f(a)$), then there exists $c \in [a, b]$ with $f(c) = L$.
\end{theorem}
\end{keyresult}

\begin{proof}[Proof of the EVT]
By Theorem~\ref{thm:cont_compact}, $f([a,b])$ is compact, hence closed
and bounded (Heine--Borel, Corollary~\ref{cor:heine_borel}). Since
$f([a,b])$ is bounded, $M = \sup f([a,b])$ exists. Since $f([a,b])$ is
closed and $M \in \cl(f([a,b]))$ (as the supremum of a bounded set is in
its closure), we have $M \in f([a,b])$. Hence there exists $x_M \in [a,b]$
with $f(x_M) = M$. The argument for the minimum is analogous, using
$m = \inf f([a,b])$.
\end{proof}

\begin{proof}[Proof of the IVT]
Assume without loss of generality that $f(a) < L < f(b)$ (if $L = f(a)$
or $L = f(b)$, take $c = a$ or $c = b$). Define
$S = \{x \in [a, b] : f(x) < L\}$. Then $S$ is nonempty (since
$a \in S$) and bounded above (by $b$). Let $c = \sup S$.

Since $f(a) < L$ and $f$ is continuous at $a$, we have $f(x) < L$ for
$x$ near $a$, so $c > a$. Similarly, $f(b) > L$ and continuity at $b$
give $c < b$. Hence $c \in (a, b)$.

Since $c = \sup S$, there exists a sequence $(x_n)$ in $S$ with
$x_n \to c$ (by the $\varepsilon$-characterization of the supremum,
Lesson~1, Remark~3.6). Since $f(x_n) < L$ and $f$ is continuous,
$f(c) = \lim f(x_n) \le L$.

If $f(c) < L$, then by continuity there exists $\delta > 0$ with
$f(x) < L$ for $x \in (c, c + \delta) \cap [a, b]$, contradicting
$c = \sup S$. Hence $f(c) = L$.
\end{proof}

\begin{rlconnection}[title={EVT in optimization and the IVT in root-finding}]
The extreme value theorem guarantees that continuous loss functions attain
their minima on compact feasible sets. This is the theoretical foundation
for constrained optimization: when we write
$\theta^* = \argmin_{\theta \in \Theta} L(\theta)$ with $\Theta$ compact
and $L$ continuous, the EVT ensures $\theta^*$ exists. The intermediate
value theorem provides fixed-point existence: if $g \colon [0,1] \to [0,1]$
is continuous, then $h(x) = g(x) - x$ satisfies $h(0) = g(0) \ge 0$ and
$h(1) = g(1) - 1 \le 0$, so the IVT gives $c$ with $h(c) = 0$, i.e.,
$g(c) = c$. This one-dimensional Brouwer fixed-point theorem underlies
existence proofs for Nash equilibria.
\end{rlconnection}

\begin{example}
\label{ex:ivt_fixed_point}
Let $g(x) = (e^x - 1)/(e - 1)$ on $[0, 1]$. Then $g(0) = 0$ and
$g(1) = 1$, and $g$ is continuous. The function $h(x) = g(x) - x$
satisfies $h(0) = 0$, so $x = 0$ is a fixed point. To find another,
note $g'(0) = 1/(e-1) < 1$, so $g(x) < x$ for small positive $x$, while
$g(1) = 1$. The IVT applied to $h$ on a suitable subinterval confirms
$x = 0$ is the unique fixed point.
\end{example}

%% ============================================================
%% SECTION 7: UNIFORM CONTINUITY
%% ============================================================
\section{Uniform Continuity}
\label{sec:uniform_cont}

\begin{definition}[Uniform Continuity]
\label{def:uniform_cont}
A function $f \colon A \to \R$ is \emph{uniformly continuous on $A$} if
for every $\varepsilon > 0$ there exists $\delta > 0$ such that
$|f(x) - f(y)| < \varepsilon$ whenever $x, y \in A$ and
$|x - y| < \delta$.
\end{definition}

\begin{remark}
\label{rem:cont_vs_uniform}
The key distinction: in ordinary continuity, $\delta$ may depend on both
$\varepsilon$ and the point $c$; in uniform continuity, $\delta$ depends
only on $\varepsilon$ and works for \emph{all} pairs of points
simultaneously.
\end{remark}

\begin{example}
\label{ex:not_uniform}
The function $f(x) = x^2$ is continuous on $\R$ but not uniformly
continuous on $\R$. Indeed, for any $\delta > 0$, choose $x = 1/\delta$
and $y = x + \delta/2$. Then $|x - y| = \delta/2 < \delta$ but
$|f(x) - f(y)| = |y^2 - x^2| = |y - x||y + x| =
(\delta/2)(2/\delta + \delta/2) \ge 1$,
so taking $\varepsilon = 1$ shows uniform continuity fails.
\end{example}

\begin{example}
\label{ex:lipschitz_uniform}
If $f$ is Lipschitz on $A$, i.e., there exists $L > 0$ with
$|f(x) - f(y)| \le L|x - y|$ for all $x, y \in A$, then $f$ is
uniformly continuous: given $\varepsilon > 0$, take $\delta =
\varepsilon / L$. For instance, $f(x) = \sin(x)$ satisfies
$|\sin(x) - \sin(y)| \le |x - y|$ (Lipschitz constant $1$), hence is
uniformly continuous on $\R$.
\end{example}

\begin{theorem}[Compact Domain Implies Uniform Continuity]
\label{thm:compact_uniform}
If $f \colon K \to \R$ is continuous and $K \subseteq \R$ is compact,
then $f$ is uniformly continuous on $K$.
\end{theorem}

\begin{proof}
Suppose for contradiction that $f$ is continuous on $K$ but not uniformly
continuous. Then there exists $\varepsilon_0 > 0$ such that for every
$n \in \N$, there exist $x_n, y_n \in K$ with $|x_n - y_n| < 1/n$ but
$|f(x_n) - f(y_n)| \ge \varepsilon_0$.

Since $K$ is compact (hence sequentially compact), $(x_n)$ has a
subsequence $x_{n_k} \to c \in K$. Since $|x_{n_k} - y_{n_k}| < 1/n_k
\to 0$, we also have $y_{n_k} \to c$.

By continuity of $f$ at $c$: $f(x_{n_k}) \to f(c)$ and
$f(y_{n_k}) \to f(c)$. Hence
$|f(x_{n_k}) - f(y_{n_k})| \to 0$,
contradicting $|f(x_{n_k}) - f(y_{n_k})| \ge \varepsilon_0 > 0$.
\end{proof}

\begin{example}
\label{ex:uniform_interval}
The function $f(x) = x^2$ is uniformly continuous on $[0, 10]$ (by
Theorem~\ref{thm:compact_uniform}, since $[0,10]$ is compact), even
though it fails to be uniformly continuous on all of $\R$
(Example~\ref{ex:not_uniform}).
\end{example}

\begin{rlconnection}[title={Uniform continuity in value function approximation}]
In RL with continuous state spaces, value functions $V^\pi(s)$ are often
shown to be Lipschitz or uniformly continuous on compact state spaces.
Uniform continuity guarantees that if we approximate $V^\pi$ on a finite
grid of states with spacing $\delta$, the approximation error is bounded
by $\varepsilon$ uniformly---there are no ``bad'' states where the error
blows up. This is essential for the correctness of discretization-based
solvers and for bounding the approximation error of neural network
function approximators on bounded domains.
\end{rlconnection}

%% ============================================================
%% SECTION 8: EXERCISES
%% ============================================================
\section{Exercises}
\label{sec:exercises}

\begin{exercise}[Classifying sets]
\textnormal{(\(*\))} \\
For each of the following subsets of $\R$, determine whether it is open,
closed, both, or neither. Find its interior, closure, and boundary.
\begin{enumerate}[nosep,label=(\alph*)]
  \item $A = [0, 1)$
  \item $B = \Z$
  \item $C = \{1/n : n \in \N\}$
  \item $D = (0, 1) \cup (1, 2)$
\end{enumerate}
\end{exercise}

\begin{exercise}[Compactness verification]
\textnormal{(\(*\))} \\
Determine which of the following sets are compact, and justify your answers
using the Heine--Borel theorem (Corollary~\ref{cor:heine_borel}):
\begin{enumerate}[nosep,label=(\alph*)]
  \item $\{x \in \R : x^2 \le 4\}$
  \item $\Q \cap [0, 1]$
  \item $\{1 - 1/n : n \in \N\} \cup \{1\}$
  \item $(0, 1]$
\end{enumerate}
\end{exercise}

\begin{exercise}[Continuity proof from the definition]
\textnormal{(\(*\))} \\
Prove directly from Definition~\ref{def:continuity} (the
$\varepsilon$-$\delta$ definition) that $f(x) = 3x + 7$ is continuous at
every $c \in \R$.
\end{exercise}

\begin{exercise}[IVT application in RL]
\textnormal{(\(**\))} \\
Let $V \colon [0, 1] \to \R$ be a continuous function with $V(0) = -1$ and
$V(1) = 5$. Let $\lambda \in (-1, 5)$ be a threshold value.
\begin{enumerate}[nosep,label=(\alph*)]
  \item Use the IVT to prove there exists $s^* \in (0, 1)$ with
    $V(s^*) = \lambda$.
  \item If $V$ is additionally assumed to be strictly increasing, prove
    that $s^*$ is unique.
\end{enumerate}
\textit{Hint:} For uniqueness, assume two distinct roots and derive a
contradiction with strict monotonicity.
\end{exercise}

\begin{exercise}[Uniform continuity failure]
\textnormal{(\(**\))} \\
Prove that $f(x) = 1/x$ is continuous but not uniformly continuous on
$(0, 1)$. Then explain why this does not contradict
Theorem~\ref{thm:compact_uniform}.
\end{exercise}

\begin{exercise}[Compactness from first principles]
\textnormal{(\(***\))} \\
Prove directly from the open-cover definition (without using Heine--Borel)
that $[0, 1]$ is compact.

\textit{Hint:} Let $\{O_\alpha\}$ be an open cover of $[0,1]$. Define
$S = \{x \in [0,1] : [0,x] \text{ is covered by finitely many }
O_\alpha\}$. Show $S$ is nonempty, $\sup S = 1$, and $1 \in S$.
\end{exercise}

%% ============================================================
%% SECTION 9: SUMMARY
%% ============================================================
\section{Summary}
\label{sec:summary}

\begin{itemize}
  \item A set $O \subseteq \R$ is \emph{open} if every point has an
    $\varepsilon$-neighborhood contained in $O$; a set is \emph{closed}
    if its complement is open, equivalently if it contains all limits of
    convergent sequences from the set.

  \item Arbitrary unions of open sets are open; finite intersections of
    open sets are open. The dual statements hold for closed sets.

  \item The \emph{closure} $\cl(A)$ consists of all limits of sequences
    in $A$; the \emph{interior} $\interior(A)$ is the largest open
    subset of $A$; the \emph{boundary} $\partial A = \cl(A) \setminus
    \interior(A)$.

  \item A set $K \subseteq \R$ is \emph{compact} if and only if it is
    closed and bounded (\emph{Heine--Borel theorem}), which is
    equivalent to sequential compactness: every sequence in $K$ has a
    subsequence converging to a point in $K$.

  \item Continuity at $c$ means $|f(x) - f(c)| < \varepsilon$ for
    $|x - c| < \delta$; equivalently, $x_n \to c$ implies
    $f(x_n) \to f(c)$. Continuous functions preserve sums, products,
    quotients, and compositions.

  \item The \emph{extreme value theorem}: a continuous function on a
    compact set attains its maximum and minimum.

  \item The \emph{intermediate value theorem}: a continuous function on
    $[a,b]$ takes every value between $f(a)$ and $f(b)$, which gives
    existence of roots and fixed points.

  \item A function is \emph{uniformly continuous} if $\delta$ depends
    only on $\varepsilon$, not on the point. Every continuous function
    on a compact set is uniformly continuous.
\end{itemize}

%% ============================================================
%% SECTION 10: FURTHER READING
%% ============================================================
\section{Further Reading}
\label{sec:further_reading}

\begin{itemize}
  \item \textbf{Lesson 3} (Differentiation and Integration) builds on
    continuity to develop differentiability, the mean value theorem,
    Taylor's theorem, and Riemann integration. The extreme value theorem
    proved here is used in the proof of Rolle's theorem, and uniform
    continuity is used to prove that continuous functions are Riemann
    integrable.

  \item In Phase~01 (Functional Analysis, Lesson~5), the concepts of
    open sets, closed sets, and compactness are generalized from $\R$ to
    arbitrary metric spaces and then to normed spaces. The Heine--Borel
    characterization fails in infinite-dimensional spaces, making
    compactness significantly harder to verify.

  \item The Arzel\`{a}--Ascoli theorem, which characterizes compact
    subsets of $C([a,b])$ (the space of continuous functions), extends
    the ideas of this lesson to function spaces and is treated in
    Phase~01.

  \item For topology in $\R^n$ and general metric spaces, see
    \citet[Chapter 2]{rudin1976}. For a more abstract topological
    perspective, see \citet[Chapter 4]{folland1999}.

  \item \citet[Chapters 3--4]{abbott2015} offers an exceptionally clear
    treatment of compactness and the extreme value theorem, with
    illuminating discussion of why each hypothesis is needed.

  \item Students proceeding to measure theory will find that open and
    closed sets generate the Borel $\sigma$-algebra $\Borel(\R)$, which
    is the natural domain for probability measures on $\R$.
\end{itemize}

\bibliography{real-analysis-refs}

\end{document}
